% ===================================================================
%  TABELUL Nr. 16 — Magazinul mare. — Bazarul. — Jucăriile.
%  Traduction 2026 d'apres texto/io, controlee sur texto/fr, texto/en
%  et texto/es. Consigne : voir texto/ro/10-tabelul-01.tex.
%
%  Le tableau d'astronomie (bloc 15-1) porte des renvois a LETTRES,
%  (a) a (f), graves sur le tableau lui-meme : ils se rangent sous le
%  numero 85, comme le dit gravuri/literi.json. Ce bloc est aussi le
%  troisieme ecart declare de renvoji.py.
%
%  DERNIER TABLEAU DE LA COLONNE ROUMAINE.
% ===================================================================

%%K t16-apar-1 apar
\VUsaut{4.0mm}
\VUtitre{15.0pt}{60}{TABELUL Nr. 16}
\VUsaut{2.0mm}
\VUpk{13.2pt}{40}{Magazinul mare. --- Bazarul.}
\VUpk{13.2pt}{40}{Jucăriile.}
\VUsaut{3.4mm}

%%K t16-01-1 p
1. --- Se apropie anul nou. Magazinele mari și-au pregătit expozițiile de
\VUgras{jucării} \textsuperscript{(1)} și de daruri de anul nou.

%%K t16-01-2 p
L-am rugat pe bunicul să mă ducă la bazar ca să văd această expoziție frumoasă, și
s-a învoit.

%%K t16-02-1 p
2. --- Mai întâi ne-am oprit înaintea frumoaselor panouri cu arme, pe care erau o
\VUgras{pușcă}, o \VUgras{șapcă}, o \VUgras{sabie}, un \VUgras{centiron} și o
pereche de \VUgras{epoleți}, sau altminteri o \VUgras{cască}, o \VUgras{platoșă}
\textsuperscript{(3)} și un \VUgras{pistol} \textsuperscript{(4)}. Toate acestea mă
interesau mult mai mult decât \VUgras{lanternele magice} \textsuperscript{(2)}.

%%K t16-tit-1 sub
\VUsaut{3.0mm}
\VUpk{10.2pt}{40}{Priveliști frumoase.}
\VUsaut{2.6mm}

%%K t16-03-1 p
3. --- În același raion stătea o \VUgras{santinelă} \textsuperscript{(5)} în
\VUgras{gheretă}; parcă păzea o menajerie întreagă de carton. Câte animale erau
acolo! \VUgras{Papagalul} \textsuperscript{(6)} pe stinghia lui,
\VUgras{rinocerul} \textsuperscript{(7)} cu cornul pe nas, \VUgras{renul}
\textsuperscript{(8)}, \VUgras{struțul} \textsuperscript{(9)}, \VUgras{maimuța}
\textsuperscript{(10)} ronțăind o nucă, \VUgras{vulpea} \textsuperscript{(11)}
ducând o găină, \VUgras{crocodilul} \textsuperscript{(12)} cu botul uriaș căscat,
\VUgras{cămila} \textsuperscript{(13)}, un \VUgras{ogar} elegant
\textsuperscript{(14)}, \VUgras{pantera} \textsuperscript{(15)}, iar apoi două
animale pe care nu le știam și care, cum mi s-a spus, sunt \VUgras{nevăstuica}
\textsuperscript{(16)} și \VUgras{hiena} \textsuperscript{(17)}. Mai era acolo și
\VUgras{leopardul} \textsuperscript{(18)} și \VUgras{lupul}
\textsuperscript{(21)}; chiar alături \VUgras{șobolani} nostimi
\textsuperscript{(19)} și \VUgras{șoricei} mărunți de cauciuc
\textsuperscript{(20)}. În sfârșit, \VUgras{hipopotamul} \textsuperscript{(22)} și
\VUgras{dromaderul} cocoșat \textsuperscript{(23)} întregeau colecția. Aș fi stat
acolo ceasuri întregi, dacă alături nu s-ar fi aflat o minunată tabără plină de
\VUgras{soldați de plumb} \textsuperscript{(29)}, care mi-a atras luarea-aminte.

%%K t16-tit-2 sub
\VUsaut{4.0mm}
\VUpk{10.2pt}{40}{Soldații și războiul.}
\VUsaut{2.8mm}

%%K t16-04-1 p
4. --- Bunicul, care este \VUgras{invalid} \textsuperscript{(24)} și a trebuit să
lase slujba din pricina unei \VUgras{răni} care i-a adus \VUgras{medalia militară},
iubește să îmi povestească \VUgras{campaniile} la care a luat parte. Mi-a lămurit
cu plăcere feluritele mișcări ale acestei mici oștiri în miniatură. Un grup de
soldați adevărați, care cercetau bazarul --- un \VUgras{genist}
\textsuperscript{(25)}, un \VUgras{vânător de munte} \textsuperscript{(26)} cu
\VUgras{bereta} lui, un \VUgras{plutonier} de infanterie \textsuperscript{(27)} și
un \VUgras{sergent} de artilerie \textsuperscript{(28)} cu \VUgras{trese} de aur pe
mânecă --- s-a oprit lângă noi ca să asculte lămuririle.

%%K t16-05-1 p
5. --- Bunicul, foarte mândru de un ascultător atât de vrednic de cinste, a
născocit pe loc un întreg plan de bătălie.

%%K t16-05-2 p
Cetatea cu \VUgras{metereze} și cu \VUgras{șanțuri} fusese împresurată de
\VUgras{oastea dușmană}, care după un lung \VUgras{bombardament} se pregătea să o
ia cu \VUgras{asalt}. Dar \VUgras{cei împresurați}, aduși la capăt de foame, au
făcut o \VUgras{ieșire} împotriva \VUgras{taberei} \VUgras{împresurătorilor}.
Respingând \VUgras{atacul} într-o \VUgras{luptă} crâncenă în jurul
\VUgras{corturilor}, împresurătorii \VUgras{biruitori} și-au trimis
\VUgras{escadroanele} în \VUgras{urmărirea} atacatorilor \VUgras{biruiți}. Aceștia
au \VUgras{dat înapoi}, lăsând \VUgras{câmpul de bătaie} presărat cu \VUgras{morți}
și cu \VUgras{răniți}. \VUgras{Sanitarii} i-au dus pe răniți la \VUgras{spitalul
de campanie}, ca să îi îngrijească \VUgras{chirurgul}.

%%K t16-05-3 p
Cu toată împotrivirea vitejească a câtorva \VUgras{batalioane} așezate în
\VUgras{careu}, \VUgras{înfrângerea} a fost deplină; ce a mai rămas din oaste a
\VUgras{fugit}, lăsând o mulțime de \VUgras{prinși}. Dușmanul se pregătea să își
încununeze \VUgras{biruința} cu luarea cetății. Poate că acesta ar fi fost
sfârșitul campaniei, și după \VUgras{armistițiu} s-ar fi putut încheia
\VUgras{tratatul de pace}.

%%K t16-tit-3 sub
\VUsaut{4.4mm}
\VUpk{10.2pt}{40}{Darurile de anul nou.}
\VUsaut{2.8mm}

%%K t16-06-1 p
6. --- Soldații tineri, care râdeau ascultând povestirea colorată a veteranului, i-au
mai pus câteva întrebări. Cât despre mine, am ocolit taraba ca să văd o
\VUgras{arbaletă} minunată \textsuperscript{(32)} și un \VUgras{arc}
\textsuperscript{(33)} cu \VUgras{săgeata} lui \textsuperscript{(31)}. Acolo am mai
găsit o colecție de animale: două \VUgras{păsări cântătoare}
\textsuperscript{(34)}, o \VUgras{privighetoare} și o \VUgras{silvie} cu resort,
care cântau cât le ținea gura, și un șir de vietăți ciudate ---
\VUgras{liliacul} \textsuperscript{(35)}, \VUgras{boa} \textsuperscript{(36)},
\VUgras{broasca țestoasă} \textsuperscript{(37)}, \VUgras{foca}
\textsuperscript{(38)}, \VUgras{balena} \textsuperscript{(39)} și
\VUgras{rechinul} \textsuperscript{(40)} din beșică.

%%K t16-07-1 p
7. --- Nu m-am uitat mult la ele, fiindcă am zărit ceea ce visam: un minunat
\VUgras{teatru de păpuși} \textsuperscript{(41)} cu \VUgras{decoruri} frumoase și cu
o \VUgras{cortină} roșie încântătoare. Pe \VUgras{scenă} doi actori parcă jucau un
\VUgras{rol} într-o \VUgras{piesă} cât se poate de interesantă, fiindcă micii
\VUgras{spectatori} de carton din loji, și de asemenea cei așezați în \VUgras{stal}
și la \VUgras{galerie} în spatele \VUgras{dirijorului}, băteau din palme cu foc.

%%K t16-07-2 p
Din păcate, acel teatru era foarte scump.

%%K t16-08-1 p
8. --- Pe deasupra era greu să alegi între jucării, fiindcă vitrinele erau pline de
jucării de tot felul: \VUgras{Arlechinul} \textsuperscript{(42)},
\VUgras{Poliținel} \textsuperscript{(43)}, \VUgras{Pierrot}
\textsuperscript{(44)}, \VUgras{bufnițe} \textsuperscript{(45)} care băteau din
aripi, \VUgras{sturzi} \textsuperscript{(46)} care fluierau, \VUgras{pudeli} care
lătrau, \VUgras{gramofonul} \textsuperscript{(47)} și \VUgras{jocurile de
răbdare}. Una dintre aceste jucării m-a înveselit mult: înfățișa o \VUgras{bătălie
navală} \textsuperscript{(48)}, în care o \VUgras{corabie} este atacată de
\VUgras{pirați} lângă un țărm sădit cu \VUgras{palmieri de cocos}.

%%K t16-noto-1 noto
\VUnotes{143.2mm}{%
\textit{(*) Vezi tabelul nr. 6.}%
}

%%K t16-09-1 p
9. --- Toate aceste minuni le puteam privi cu totul în tihnă, fiindcă în magazin era
puțină lume --- abia un pumn de gură-cască, un \VUgras{dragon}
\textsuperscript{(49)} și un \VUgras{încasator} \textsuperscript{(50)}, care
înfățișa o \VUgras{poliță} la \VUgras{casa} principală \textsuperscript{(51)}.

%%K t16-10-1 p
10. --- L-am întrebat pe bunicul la ce sunt cărțile de pe rafturile din spatele
\VUgras{casierului}; mi-a spus că sunt \VUgras{registrele de contabilitate}.

%%K t16-tit-4 sub
\VUsaut{3.6mm}
\VUpk{10.2pt}{40}{Cercetând magazinul.}
\VUsaut{2.8mm}

%%K t16-11-1 p
11. --- Ieșind din raionul de jucării, am mers la cel de șelărie, ca să vedem
\VUgras{șeile} \textsuperscript{(53)}, \VUgras{scările}
\textsuperscript{(54)}, \VUgras{frâiele} \textsuperscript{(55)},
\VUgras{zăbalele} \textsuperscript{(56)} și \VUgras{cravașele}. În vreme ce
\VUgras{șeful de raion} \textsuperscript{(57)} îi dădea bunicului lămuriri, am mers
să văd expoziția de \VUgras{porțelan} și de \VUgras{sticlă}
\textsuperscript{(58)}, unde erau \VUgras{salatiere} frumoase și \VUgras{ceainice}
elegante \textsuperscript{(59)}.

%%K t16-12-1 p
12. --- Ca să urcăm la etajul întâi, am trecut prin spatele vitrinei de
\VUgras{corsete} \textsuperscript{(60)}, lângă raionul de ciorapi, unde erau
așezați \VUgras{pantaloni} frumoși \textsuperscript{(63)}. Alături o
\VUgras{vânzătoare} \textsuperscript{(61)} o ajuta pe o \VUgras{cumpărătoare}
\textsuperscript{(62)} să aleagă \VUgras{stofe} frumoase de mătase
\textsuperscript{(64)}.

%%K t16-12-2 p
Urcând pe scară, am băgat de seamă că magazinul este împodobit cu
\VUgras{lampioane} japoneze \textsuperscript{(65)}, cu \VUgras{umbreluțe}
\textsuperscript{(66)} și cu \VUgras{palmieri} uriași \textsuperscript{(70)}.

%%K t16-13-1 p
13. --- La etajul întâi nu prea era ce vedea: numai mobilă (*), \VUgras{case de
bani} \textsuperscript{(67)}, \VUgras{scrinuri} \textsuperscript{(68)} și
\VUgras{cărucioare de copii} \textsuperscript{(69)}. În schimb, priveliștea de
obște a magazinului era foarte \VUgras{plăcută}.

%%K t16-14-1 p
14. --- Sub noi vedeam instrumentele muzicale: \VUgras{harpe}
\textsuperscript{(71)}, \VUgras{chitare} \textsuperscript{(72)},
\VUgras{mandoline} \textsuperscript{(73)}, \VUgras{cornete}
\textsuperscript{(74)}, \VUgras{clarinete} \textsuperscript{(75)}, viori cu
\VUgras{arcușurile} lor \textsuperscript{(76)} și chiar \VUgras{toba mare}
\textsuperscript{(77)}. Ceva mai departe, la taraba de papetărie, un
\VUgras{student} \textsuperscript{(78)} se tocmea cu \VUgras{vânzătorul}
\textsuperscript{(79)} pentru o mapă de hârtii. Lângă ei deosebeam un
\VUgras{abecedar} frumos \textsuperscript{(80)}.

%%K t16-14-2 p
În mijlocul magazinului stătea un \VUgras{pom de Crăciun} înalt
\textsuperscript{(81)}, tot încărcat cu lumânări, cu podoabe și cu jucării de tot
felul: \VUgras{căluți de lemn} \textsuperscript{(82)}, \VUgras{păpuși}
\textsuperscript{(83)} și altele.

%%K t16-14-3 p
\VUgras{Raionul de parfumerie} \textsuperscript{(84)} cu \VUgras{apele de dinți} și
cu feluritele \VUgras{parfumuri} răspândea un miros foarte plăcut, care se ridica
până la noi.

%%K t16-tit-5 sub
\VUsaut{3.4mm}
\VUpk{10.2pt}{40}{Cumpărăm ceva.}
\VUsaut{2.8mm}

%%K t16-15-1 p
15. --- Fiindcă bunicului începea să îi fie urât, am coborât din nou pe scara cea
mare. S-a oprit o clipă înaintea \VUgras{tabelului de astronomie}
\textsuperscript{(85)}, pe care, cum mi-a spus, erau înfățișate \VUgras{cometele}
\textsuperscript{(a)}, \VUgras{eclipsele de lună și de soare}
\textsuperscript{(b)}, roza vânturilor cu \VUgras{cele patru zări}
\textsuperscript{(c)} \VUgras{(nordul, sudul, răsăritul și apusul)}, harta celor
două \VUgras{emisfere} \textsuperscript{(d)}, \VUgras{planetele}
\textsuperscript{(e)} și \VUgras{sistemul solar} \textsuperscript{(f)}. Eu nu am
înțeles nimic din toate acestea, și mă interesa mult mai mult livreaua cu galoane a
\VUgras{comisionarului} \textsuperscript{(86)} care trecea pe lângă noi și
\VUgras{evantaiele} drăguțe \textsuperscript{(87)} de lângă mine, așezate lângă
\VUgras{portmonee} \textsuperscript{(88)} și \VUgras{portofele}
\textsuperscript{(89)}.

%%K t16-16-1 p
16. --- Am admirat de asemenea pe tarabă \VUgras{penele} \textsuperscript{(90)} și
\VUgras{cataramele de curea} \textsuperscript{(91)}, și \VUgras{florile
artificiale} drăguțe \textsuperscript{(94)}, așezate cu gust în coșuri.
\VUgras{Violetele}, \VUgras{micșunelele}, \VUgras{zambilele}
\textsuperscript{(95)}, \VUgras{mușcatele} \textsuperscript{(96)},
\VUgras{crinii} \textsuperscript{(97)} și \VUgras{condurii doamnei}
\textsuperscript{(98)} erau foarte bine imitate.

%%K t16-tit-6 sub
\VUsaut{4.4mm}
\VUpk{10.2pt}{40}{Darul bunicului.}
\VUsaut{2.8mm}

%%K t16-17-1 p
17. --- L-am rugat pe bunicul să îmi cumpere una dintre \VUgras{periuțele de dinți}
\textsuperscript{(92)} din cutia de lângă \VUgras{ace cu gămălie}
\textsuperscript{(93)}. S-a învoit, și am mers eu însumi să plătesc cumpărătura la
casă, lângă care se lega un \VUgras{pachet} mare \textsuperscript{(100)} pentru un
\VUgras{cumpărător} \textsuperscript{(99)}.

%%K t16-18-1 p
18. --- Deodată, printre oamenii care stăteau lângă \VUgras{bronzurile}
\textsuperscript{(101)} de artă, s-a iscat o ușoară tulburare. Un \VUgras{hoț}
\textsuperscript{(102)} tocmai fusese prins asupra faptului de \VUgras{detectivul
magazinului} \textsuperscript{(103)} în clipa când încerca să fure de la cineva. Au
trimis după sergent ca să îl \VUgras{rețină}, și tocmai când ieșeam, vinovatul era
dus la secție.
\VUsaut{6.0mm}
\VUfilet{22mm}
